\section{Shooting through the coupling bands}
\label{sec:surviving-interval-shooting}
A nonlinear step is available when its input coupling belongs to the
prescribed band. Its output need not belong to the next band for every
input parameter. We construct a surviving interval at each level; this
preserves the activity continuity needed for the subsequent step.

Fix the scale sequence and all the preceding nonlinear choices. Let
$s=s_G>0$, $L_i\ge2$ be the integer block sides, and set
\[
 S_0=0,\qquad S_{i+1}=S_i+\log L_i,\qquad
 \widehat\beta_i=\bar\beta+sS_i,\qquad \bar\beta>0.
\]
The permitted band is
$[\widehat\beta_i/2,2\widehat\beta_i]$.
Suppose the purchased initial family exists for every cutoff and period,
and that the continuous return of the preceding section applies on
any compact interval whose inputs satisfy the history conditions and
belong to the fine band. Preserve its exact density identity and its
retained next-step conditions. The adjacent coupling bound has the form
\begin{equation}\label{eq:shoot-adjacent-drift}
 |\beta_{i+1}-\beta_i-s\log L_i|\le\varepsilon_{i+1},
 \qquad \varepsilon_{i+1}\le\tfrac{s}{4}\log L_i.
\end{equation}
Here $\varepsilon_{i+1}\ge0$ is the drift allowance of the stage
which the fixed schedule assigns to level $i+1$; it is a stage
constant, chosen before the block lengths. The first inequality is a
bound on the actual returned representations: it combines the stage
bound for the scalar shift $\alpha_{N,i}$ with the purchased reset
width of the preceding section. The second inequality is a property
of the schedule, whose block sides are chosen after the stage drift
allowances.

\subsection{Restricting a curve without crossing a barrier}
We first record the elementary interval argument. Let $f$ be continuous
on $[x,y]$, and suppose
$f(x)\le d\le u\le f(y)$. There are $z\le w$ in $[x,y]$ such that
\[
 f(z)=d,\quad f(w)=u,\quad f([z,w])\subset[d,u].
\]
The case $d=u$ follows by taking one point of the nonempty level set.
Otherwise let $w$ be the first point of $[x,y]$ at which $f=u$.
The level set is nonempty by the intermediate value theorem, and
compactness makes its minimum belong to it. Let $z$ be the last point
of $[x,w]$ at which $f=d$; this level set is likewise nonempty and
compact. If $f(c)<d$ at some $c\in[z,w]$, the intermediate value
theorem between $c$ and $w$ gives a later point with value $d$,
contrary to the choice of $z$. If $f(c)>u$, the same theorem between
$x$ and $c$ gives an earlier point with value $u$, contrary to the
choice of $w$. This proves the claim without any monotonicity assumption
on $f$.

Define inner barriers
\[
 \ell_i=\bar\beta+\tfrac{s}{2}S_i,
 \qquad u_i=\bar\beta+2sS_i.
\]
Both lie in the permitted band, since $S_i\ge0$. If an input at
level $i+1$ has coupling $\ell_{i+1}$, its returned coarse coupling
satisfies, by \eqref{eq:shoot-adjacent-drift},
\[
 \beta_i\le\ell_{i+1}-s\log L_i+\varepsilon_{i+1}
             \le\ell_i-\tfrac{s}{4}\log L_i<\ell_i.
\]
If the input coupling is $u_{i+1}$, the corresponding lower bound is
\[
 \beta_i\ge u_{i+1}-s\log L_i-\varepsilon_{i+1}
             \ge u_i+\tfrac{3s}{4}\log L_i>u_i.
\]
Thus the two endpoints cross the next two barriers in the required
order, even though the intermediate output curve may oscillate.

\subsection{Descending through a finite history}
At cutoff $N$ restrict the initial bare-action family to
$[\ell_N,u_N]$. Its coupling is its parameter, it is continuous in the
labelled activity norms, and it satisfies the prescribed initial
conditions. Suppose a compact parameter interval already carries a
family of histories from $N$ down to $i+1$, with every represented
coupling in its band, and with the endpoint couplings at level $i+1$
equal to $\ell_{i+1},u_{i+1}$. Apply the continuous nonlinear step to
the level-$i+1$ curve on this interval. The endpoint inequalities just
proved and the interval lemma give a subinterval on which the output
coupling stays between $\ell_i,u_i$, attaining these values at the
two ends.

Restrict every earlier representation in the history to this
subinterval and append the newly constructed level-$i$ representation.
Restriction preserves the fixed dictionaries, their norm continuity,
all activity bounds and every previously established adjacent identity.
The new step supplies the same properties at level $i$. Induction
therefore reaches level zero. Since $\ell_0=u_0=\bar\beta$, either
endpoint supplies a history with exactly the desired coarse coupling.
No representation was extended across an interval on which its input
band condition failed.

There is also one fixed neighborhood of coarse targets that uses the
same selected data. Put
\[
 h=\min\{\bar\beta/2,s\log L_0/4\}>0.
\]
For any $|b-\bar\beta|\le h$, descend only to level $1$ and apply the
last nonlinear step there. Its lower endpoint returns at most
$\bar\beta-s\log L_0/4\le b$, and its upper endpoint returns at
least $\bar\beta+3s\log L_0/4\ge b$. The intermediate value theorem
selects an input parameter with output coupling $b$. That output lies
in the level-zero band because $|b-\bar\beta|\le\bar\beta/2$.
For $N=0$, use the initial family directly at parameter $b$.
This construction asserts existence, not uniqueness of a shooting
history or continuity of its choice as $b$ varies.

For every cutoff and period select one such history. The selected
paths depend on cutoff and level, and were fixed before the period;
therefore their identity across periods is preserved. Each selected
history starts at the bare Wilson density with vacuum parameter zero,
satisfies all its local activity bounds, and has the exact full-density
transfer at each step. The stronger retained conditions are preserved
by taking them as part of the next-step requirement throughout the
interval induction. The induction selects no new bridges, paths or
support exponents.

\subsection{The running-coupling estimate}
Let
$E_0=0$ and $E_i=\sum_{j<i}\varepsilon_{j+1}$.
Telescoping \eqref{eq:shoot-adjacent-drift} gives, uniformly over every
selected cutoff $N\ge i$ and period,
\begin{equation}\label{eq:shoot-running-error}
            |\beta_i-\widehat\beta_i|\le E_i+|b-\bar\beta|.
\end{equation}
Let $J_i$ denote the stage assigned to level $i$ by the fixed
schedule. The schedule has two further properties, which are inputs
here: the stages are unbounded along the levels, and the block side
of each stage is chosen so large, after its drift allowance, that
\begin{equation}\label{eq:shoot-drift-ratio}
 \frac{\varepsilon_{i+1}}{\log L_i}\le J_{i+1}^{-2},
                  \qquad J_i\longrightarrow\infty.
\end{equation}
The allowances $\varepsilon_{i+1}$ are the same as in
\eqref{eq:shoot-adjacent-drift}.
Fix $\epsilon>0$ and choose $j$ so that $J_{i+1}^{-2}\le\epsilon/2$
for $i\ge j$. Then
\[
 E_i\le E_j+\tfrac\epsilon2(S_i-S_j),\qquad i\ge j.
\]
Since $S_i\to\infty$, eventually
$E_j+|b-\bar\beta|\le\epsilon S_i/2$. Substitution into
\eqref{eq:shoot-running-error} yields
\[
                 |\beta_i-\widehat\beta_i|\le\epsilon S_i
\]
uniformly in cutoff and period. In particular the error is $o(S_i)$.
There is no assertion that the sum of the step errors is uniformly
bounded in the depth. The estimate uses their small ratio to the
accumulated logarithmic scale.
